% =============================================================================
%  Seção IV — Resultados e Discussão
%  Artigo: Detector de Bordas Canny com Seleção Adaptativa de Parâmetros em FPGA
%  Referência: Kalbasi & Nikmehr, IEEE Access, 2020
%
%  ATENÇÃO — pré-requisito no preamble do main.tex:
%    \usepackage[caption=false,font=footnotesize]{subfig}
%  (remover \usepackage{subcaption} se existir — incompatível com IEEEtran)
% =============================================================================

\section{Resultados e Discussão}
\label{sec:resultados}

Esta seção apresenta a validação experimental da arquitetura implementada em RTL Verilog,
organizada em quatro vertentes: ambiente de simulação, validação funcional do pipeline,
eficácia da Seleção Adaptativa de Parâmetros (APS) e desempenho temporal.

% ---------------------------------------------------------------------------
\subsection{Ambiente de Simulação e Conjunto de Imagens}
\label{ssec:ambiente}
% ---------------------------------------------------------------------------

Toda a verificação foi conduzida em simulação comportamental com o simulador iverilog 12.0
(\textit{open-source}, sem licença), que executa o RTL sintetizável descrito em Verilog-2012.
A ferramenta de síntese lógica disponível no laboratório é o Cadence Genus (ambiente acadêmico);
porém, os tempos de acesso ao ambiente Cadence inviabilizaram a síntese completa dentro do
cronograma deste trabalho. As métricas de desempenho apresentadas nas próximas subseções
são, portanto, projeções baseadas em contagem de ciclos extraída da simulação RTL.

O conjunto de imagens de avaliação é composto por seis fotografias coloridas de domínio
público amplamente utilizadas em literatura de processamento de imagens:
\textit{Lena}, \textit{Peppers}, \textit{Stones}, \textit{Bee}, \textit{Brick Wall} e
\textit{Ducks}. Todas foram redimensionadas para $512 \times 512$ pixels.
A conversão de cor (RGB~$\to$~cinza) é realizada pelo próprio \textit{datapath} de
pré-processamento da arquitetura (módulo \texttt{rgb2gray.v}), replicando a fórmula
ITU-R~BT.601 via soma ponderada com operadores de deslocamento.

A Figura~\ref{fig:bordas_exemplos} exibe quatro pares (imagem original / mapa de bordas RTL)
gerados no modo APS ($\mathrm{MDP}=94\%$), ilustrando a capacidade do \textit{pipeline}
de extrair estruturas de borda em diferentes conteúdos visuais.

% figure* ocupa as duas colunas do formato IEEE — só pode ser posicionada no topo (t)
\begin{figure*}[t!]
    \centering
    % Linha 1: Lena e Bee
    \subfloat[\textit{Lena} — original\label{fig:lena_orig}]{%
        \includegraphics[width=0.23\linewidth]{secoes/resultados/lena_original.png}}\hfill
    \subfloat[\textit{Lena} — bordas RTL\label{fig:lena_edges}]{%
        \includegraphics[width=0.23\linewidth]{secoes/resultados/lena_edges.png}}\hfill
    \subfloat[\textit{Bee} — original\label{fig:bee_orig}]{%
        \includegraphics[width=0.23\linewidth]{secoes/resultados/bee_original.png}}\hfill
    \subfloat[\textit{Bee} — bordas RTL\label{fig:bee_edges}]{%
        \includegraphics[width=0.23\linewidth]{secoes/resultados/bee_edges.png}}\\[4pt]
    % Linha 2: Peppers e Brick Wall
    \subfloat[\textit{Peppers} — original\label{fig:peppers_orig}]{%
        \includegraphics[width=0.23\linewidth]{secoes/resultados/peppers_original.png}}\hfill
    \subfloat[\textit{Peppers} — bordas RTL\label{fig:peppers_edges}]{%
        \includegraphics[width=0.23\linewidth]{secoes/resultados/peppers_edge.png}}\hfill
    \subfloat[\textit{Brick Wall} — original\label{fig:wall_orig}]{%
        \includegraphics[width=0.23\linewidth]{secoes/resultados/wall_original.png}}\hfill
    \subfloat[\textit{Brick Wall} — bordas RTL\label{fig:wall_edges}]{%
        \includegraphics[width=0.23\linewidth]{secoes/resultados/wall_edge.png}}
    \caption{Pares imagem original / mapa de bordas RTL para quatro imagens do conjunto
             de avaliação (modo APS, $\mathrm{MDP}=94\%$, $512\times512$~px).}
    \label{fig:bordas_exemplos}
\end{figure*}

% ---------------------------------------------------------------------------
\subsection{Validação Funcional do Pipeline}
\label{ssec:validacao}
% ---------------------------------------------------------------------------

A Figura~\ref{fig:validacao_funcional} compara, para a imagem \textit{Lena} em
$512\times512$~px, a entrada em escala de cinza, a referência gerada pelo OpenCV
\texttt{cv2.Canny} e a saída da simulação RTL operando no modo APS com
$\mathrm{MDP}=94\%$.

\begin{figure}[htbp]
    \centering
    \subfloat[Entrada (cinza)\label{fig:val_entrada}]{%
        \includegraphics[width=0.31\linewidth]{secoes/resultados/lena_original.png}}\hfill
    \subfloat[Ref.\ OpenCV ($T_h=38$, $T_l=15$)\label{fig:val_opencv}]{%
        \includegraphics[width=0.31\linewidth]{secoes/resultados/lena_opencv.jpg}}\hfill
    \subfloat[Saída RTL (APS, MDP=$94\%$)\label{fig:val_rtl}]{%
        \includegraphics[width=0.31\linewidth]{secoes/resultados/lena_edges.png}}
    \caption{Validação funcional: entrada em escala de cinza, referência OpenCV e
             saída da simulação RTL. Imagem \textit{Lena}, $512\times512$~px.}
    \label{fig:validacao_funcional}
\end{figure}

Qualitativamente, ambas as saídas capturam as principais estruturas de borda da imagem:
o contorno do chapéu, o ombro e os traços faciais. A saída RTL exibe menor densidade de
bordas detectadas (12.842~pixels, $\approx 5\%$) em relação à referência OpenCV a limiares
equivalentes (22.664~pixels, $\approx 8{,}8\%$), o que é esperado, pois as implementações
de Supressão de Não-Máximos (NMS) e histerese diferem entre o \textit{hardware} descrito
e a biblioteca de referência.

A métrica de sobreposição (IoU) entre os dois mapas de borda, calculada sobre a região
de saída válida do pipeline ($(W{-}6)\times(H{-}5)$ pixels), foi de $0{,}0595$, valor
que reflete as diferenças estruturais entre as duas implementações e não caracteriza erro
funcional. É importante notar que o detector de bordas Canny é intrinsecamente
sensível aos parâmetros: pequenas variações nos limiares ou no operador de gradiente
produzem diferenças significativas nos mapas de borda, tornando a comparação direta
de pixels pouco informativa sem alinhamento de parâmetros.

% ---------------------------------------------------------------------------
\subsection{Eficácia da Seleção Adaptativa de Parâmetros}
\label{ssec:eficacia}
% ---------------------------------------------------------------------------

Para avaliar o benefício da APS em condições de ruído variável, foram realizadas
28~simulações RTL: 14 níveis de ruído sal-e-pimenta ($5\%$ a $70\%$, passo de $5\%$)
$\times$ 2 modos (limiares fixos vs.\ APS). O ruído sal-e-pimenta é um modelo
impulsivo no qual pixels aleatórios são corrompidos para os valores extremos da faixa
de intensidade (0 ou 255). Esse modelo é o adequado para este sistema, pois o estimador
de ruído do hardware (\texttt{noise\_estimator.v}) opera contando pixels com valores
próximos aos extremos — exatamente a assinatura do ruído sal-e-pimenta.

Em cada rodada, a imagem \textit{Lena} ($128\times128$~px) teve ruído injetado via
\textit{script} Python, sendo convertida para hexadecimal e injetada no \textit{testbench}
RTL sem recompilação do projeto.

Os parâmetros comparados são:

\begin{itemize}
    \item \textbf{Fixo}: $T_h = 64$, $T_l = 26$ (correspondentes a $\mathrm{MDP}=91\%$,
          ruído$=5\%$ na tabela APS), mantidos constantes para todos os níveis de ruído.
    \item \textbf{APS}: limiares selecionados automaticamente pela tabela
          pré-calculada (Tabela~\ref{tab:fom_resultados}), com $\mathrm{MDP}=94\%$.
          O estimador de ruído integrado ao \textit{hardware} detecta o nível de
          contaminação e indexa a entrada adequada da LUT.
\end{itemize}

A qualidade de detecção é quantificada pela \textit{Figure of Merit} de
Pratt~\cite{pratt1978digital}:
\begin{equation}
    \mathrm{FOM} = \frac{1}{\max(|I|,\,|R|)}
                   \sum_{i \in R} \frac{1}{1 + \alpha\, d^2(i)},
    \label{eq:fom}
\end{equation}
onde $|I|$ é o número de pixels de borda no gabarito, $|R|$ é o número detectado,
$d(i)$ é a distância euclidiana do pixel detectado $i$ ao pixel de borda mais próximo
no gabarito, e $\alpha = 1/9$. O gabarito é gerado pelo \texttt{cv2.Canny} na imagem
limpa com $T_h = 300$, $T_l = 120$ ($\approx 3{,}8\%$ de densidade de bordas).

\begin{table}[htbp]
    \centering
    \caption{FOM de Pratt ($\alpha=1/9$): Parâmetros Fixos vs.\ APS ($\mathrm{MDP}=94\%$).
             Imagem: \textit{Lena} $128\times128$~px, ruído sal-e-pimenta.}
    \label{tab:fom_resultados}
    \renewcommand{\arraystretch}{1.12}
    \setlength{\tabcolsep}{6pt}
    \begin{tabular}{cccccc}
        \toprule
        \textbf{Ruído} & \textbf{$\sigma_\mathrm{APS}$} & \textbf{$T_h^\mathrm{APS}$}
            & \textbf{FOM Fixo} & \textbf{FOM APS} & \textbf{Ganho} \\
        \midrule
         5\% & 1{,}3 &  38 & 0{,}2555 & 0{,}2484 & $-$0{,}0071 \\
        10\% & 1{,}3 &  49 & 0{,}2477 & 0{,}2349 & $-$0{,}0128 \\
        15\% & 1{,}3 &  54 & 0{,}2275 & 0{,}2223 & $-$0{,}0052 \\
        20\% & 1{,}5 &  49 & 0{,}2354 & 0{,}2332 & $-$0{,}0022 \\
        25\% & 1{,}5 &  49 & 0{,}2316 & 0{,}2341 & $+$0{,}0025 \\
        30\% & 1{,}5 &  64 & 0{,}2331 & 0{,}2322 & $-$0{,}0010 \\
        35\% & 1{,}5 &  64 & 0{,}2386 & 0{,}2384 & $-$0{,}0002 \\
        40\% & 1{,}5 &  64 & 0{,}2286 & 0{,}2265 & $-$0{,}0021 \\
        45\% & 1{,}5 &  74 & 0{,}2281 & 0{,}2299 & $+$0{,}0018 \\
        50\% & 1{,}5 &  74 & 0{,}2291 & 0{,}2316 & $+$0{,}0025 \\
        55\% & 1{,}5 &  84 & 0{,}2365 & 0{,}2376 & $+$0{,}0011 \\
        60\% & 1{,}5 &  95 & 0{,}2296 & 0{,}2339 & $+$0{,}0043 \\
        65\% & 1{,}5 & 105 & 0{,}2402 & 0{,}2404 & $+$0{,}0002 \\
        70\% & 1{,}5 & 100 & 0{,}2316 & 0{,}2359 & $+$0{,}0043 \\
        \bottomrule
    \end{tabular}
\end{table}

Os resultados da Tabela~\ref{tab:fom_resultados} revelam duas regiões de comportamento
distintas: em baixo ruído ($\leq 20\%$), os limiares fixos superaram ligeiramente o modo
APS, porque $T_h=64$ gera um mapa mais esparso, mais próximo do gabarito de alta limiar.
Em ruído moderado a severo ($\geq 25\%$), o modo APS supera os limiares fixos de forma
consistente, com ganho acumulado médio de $+0{,}0024$ na faixa $25$--$70\%$. A APS eleva
progressivamente $T_h$ de 38 (para 5\% de ruído) até 105 (para 65\%), filtrando falsos
positivos induzidos pelo ruído sal-e-pimenta.

\begin{figure}[htbp]
    \centering
    \includegraphics[width=\linewidth]{secoes/resultados/aps_efficacy.png}
    \caption{FOM de Pratt ($\alpha=1/9$) em função do nível de ruído sal-e-pimenta
             para parâmetros fixos ($T_h=64$) e APS ($\mathrm{MDP}=94\%$).
             Imagem \textit{Lena}, $128\times128$~px.}
    \label{fig:aps_efficacy}
\end{figure}

% ---------------------------------------------------------------------------
\subsection{Análise de Desempenho Temporal}
\label{ssec:desempenho}
% ---------------------------------------------------------------------------

A latência do pipeline foi medida instrumentando o \textit{testbench} com um contador
de ciclos de clock (\texttt{cycle\_count}), exibido ao final da simulação via
\texttt{\$display("CICLOS\_TOTAIS:\%0d", cycle\_count)}. A simulação de referência
utilizou a imagem \textit{Lena} em $64\times64$~px e retornou:

\begin{equation*}
    N_\text{ciclos}^{64\times64} = 236.716\ \text{ciclos de clock.}
\end{equation*}

Dividindo pelo número de pixels injetados ($64\times64 = 4.096$) obtém-se uma cadência
de $\approx 57{,}8$~ciclos/pixel, com o gargalo localizado na operação de convolução
gaussiana $7\times7$ (49~ciclos de MAC mais 2~transições de estado da FSM de suavização).

Extrapolando linearmente para uma imagem $512\times512$ ($262.144$~pixels):

\begin{equation*}
    N_\text{ciclos}^{512\times512} \approx 57{,}8 \times 262.144 \approx 15{,}15 \times 10^6\ \text{ciclos.}
\end{equation*}

A Tabela~\ref{tab:desempenho} consolida a projeção de FPS para três frequências de
operação representativas:

\begin{table}[htbp]
    \centering
    \caption{Projeção de quadros por segundo (FPS) para imagens $512\times512$
             em função da frequência de operação. Valores baseados em contagem de
             ciclos da simulação RTL (não em medição em hardware real).}
    \label{tab:desempenho}
    \renewcommand{\arraystretch}{1.12}
    \begin{tabular}{ccc}
        \toprule
        \textbf{Frequência (MHz)} & \textbf{Ciclos} & \textbf{FPS (projetado)} \\
        \midrule
        100 & $15{,}15 \times 10^6$ & $6{,}6$ \\
        150 & $15{,}15 \times 10^6$ & $9{,}9$ \\
        200 & $15{,}15 \times 10^6$ & $13{,}2$ \\
        \bottomrule
    \end{tabular}
\end{table}

O resultado de $6{,}6$~FPS a $100$~MHz posiciona a arquitetura na faixa de aplicações
de inspeção industrial de baixa velocidade e robótica embarcada. Para atingir taxas de
vídeo padrão ($30$~FPS a $512\times512$), a frequência de operação precisaria ser de
aproximadamente $455$~MHz, o que demanda um dispositivo moderno (ex.\ Xilinx UltraScale+
ou Intel Agilex), cuja $F_\text{max}$ pode ser estimada após síntese e \textit{place-and-route}.

Cabe ressaltar que o \textit{pipeline} processa um pixel por ciclo de clock assim que o
pipeline está aquecido (\textit{throughput} de 1 px/ciclo após $\approx 57$ ciclos de
latência inicial). O gargalo de desempenho reside exclusivamente no filtro gaussiano
$7\times7$, cujo datapath serial processa os 49 coeficientes sequencialmente.
Uma versão paralela do filtro eliminaria esse gargalo ao custo de maior área em LUTs.
